\documentclass[11pt,a4paper]{article}
\usepackage[utf8]{inputenc}
\usepackage[T1]{fontenc}
\usepackage{lmodern,amsmath,amssymb,textcomp}
\usepackage[margin=24mm]{geometry}
\usepackage{longtable,booktabs,array,enumitem}
\usepackage[hidelinks]{hyperref}
\setlength{\parindent}{0pt}
\setlength{\parskip}{6pt}
\emergencystretch=3em
\begin{document}
{\LARGE\bfseries Finite order obstructions for edge-resilient four-critical graphs\par}\medskip
{\small Provisional mathematical authors: Rachel Teo, Wenyi Wang, Hamdi Abderrahmene\par}\medskip
{\small Judging and evaluation record: the submissions discussed in this document have been judged and evaluated by Alejandro Zarzuelo Urdiales for OpenMath 2026. This dated review records the stated verification, classification and unresolved holds; it is a preliminary evaluation rather than a signed award decision.\par}

{\small Event provenance: OpenMath 2026 ran 27 September-2 October 2026 with worldwide online participation. Detailed organizer listings specify Harvard Innovation Labs for the 27 September in-person event; the OpenMath programme advertises a hybrid kickoff at MIT. Sundai identifies its Harvard/MIT community. Organizer sources: \url{https://rsihouse.ai/openmath} ; \url{https://luma.com/yzp9abvr} ; \url{https://www.sundai.club/events/boston/recursive-learning-hack-with-harvard-innovation-labs}\par}

{\small Rachel Teo  -  Sakana AI.\par}

{\small Wenyi Wang  -  Sakana AI.\par}

{\small Hamdi Abderrahmene  -  Sakana AI.\par}

{\small Affiliations follow the recorded author declarations or public institutional/profile sources. Preferred author names, author order and publication affiliations await author review. This editorial compilation does not establish anthology coauthorship.\par}

{\small Original-result working manuscript | OpenMath 2026 | 5 October 2026\par}\medskip
{\small Central source and adjudication archive: \url{https://github.com/alejandrozu/openmath-2026-judging}\par}

\subsection{Abstract}
A finite vertex-critical four-chromatic graph that stays four-chromatic after deletion of any r or fewer edges satisfies 27r+19\ensuremath{\leq}4n, together with stronger quadratic and integral independent-set filters. These refine explicit finite lower bounds, without improving the prior eventual asymptotic restriction.

\subsection{Graph hypotheses and inequalities}
Let G be a finite simple graph on n vertices, with chromatic number four. Assume that deletion of any vertex lowers that number, and that deletion of any set of at most r edges does not lower it, where r\ensuremath{\geq}1. The submitted theorem gives 27r+19\ensuremath{\leq}4n.

Two further necessary conditions are 27(r+1)(n\ensuremath{-}1)\ensuremath{\leq}(2n+1)\ensuremath{^2} and 9b(r+1)\ensuremath{\leq}(2b+1)(n\ensuremath{-}b), where b=floor((n+1)/3). The integral filter can exclude an order allowed by the linear inequality: at r=4 it requires n\ensuremath{\geq}33, whereas the linear inequality alone gives n\ensuremath{\geq}32. These are one family of obstructions, not three separate original results.

\subsection{Criticality supplies a finite code}
Deleting a vertex of a four-critical graph produces a three-colouring of the remainder. The proof extracts a large independent colour class and encodes neighbourhood information on that class. Edge resilience forces separation between the relevant codewords: too few differences would produce a critical edge set of size at most r.

The finite-code theorem compares a lower bound for total pairwise distance with an upper bound obtained by coordinate counts. This is a Plotkin/Cauchy-Schwarz argument. Optimizing the independent-set parameter and retaining its integer size gives the quadratic and exact filters; a final numerical derivation yields the clean linear inequality.

The formal graph module supplies every hypothesis of the numeric code lemma from the actual canonical criticality definitions. The final theorem is not a numeric inequality assuming an unproved graph-to-code model. The target audit checks deletion and finite-cardinality conventions against the registered problem.

\subsection{An ordinary derivation of the code and order bounds}
Put s=r+1\ensuremath{\geq}2. For any vertex v, choose a proper three-colouring of G\ensuremath{-}v. Each of its three colour classes contains at least s neighbours of v: otherwise delete the fewer than s edges from v into that class and give v its colour, producing a three-colouring after too few deletions. Thus every vertex has degree at least 3s.

Choose a maximum independent set A of size a and put H=V(G)\ensuremath{\setminus}A, |H|=h=n\ensuremath{-}a. A three-colouring after deleting one vertex shows n\ensuremath{-}1\ensuremath{\leq}3a. For each v\ensuremath{\in}A let C\_v\ensuremath{\subseteq}H be its nonneighbours in H. Independence of A and the degree bound imply |C\_v|\ensuremath{\leq}h\ensuremath{-}3s. If u,v\ensuremath{\in}A are distinct, take a three-colouring of G\ensuremath{-}u and look at neighbours x of u having the colour of v. There are at least s of them. Such an x is outside A and cannot be adjacent to v, so it belongs to C\_v\ensuremath{\setminus}C\_u. Therefore |C\_v\ensuremath{\setminus}C\_u|\ensuremath{\geq}s in both directed orders.

Here is the complete numerical lemma. Let a subsets C\_v of an h-element ground set satisfy those weight and directed separation bounds. Write t\_x=\#\{v:x\ensuremath{\in}C\_v\}, T=\ensuremath{\sum}t\_x and Q=\ensuremath{\sum}t\_x\ensuremath{^2}. Double counting directed differences gives a(a\ensuremath{-}1)s\ensuremath{\leq}\ensuremath{\sum}\_(u,v)|C\_v\ensuremath{\setminus}C\_u|=aT\ensuremath{-}Q. Also T\ensuremath{\leq}a(h\ensuremath{-}3s) and T\ensuremath{^2}\ensuremath{\leq}hQ by Cauchy-Schwarz.

If h\ensuremath{\leq}6s, then T\ensuremath{\leq}a(h\ensuremath{-}3s)\ensuremath{\leq}ah/2. The function aT\ensuremath{-}T\ensuremath{^2}/h is increasing for 0\ensuremath{\leq}T\ensuremath{\leq}ah/2, so a(a\ensuremath{-}1)s\ensuremath{\leq}a\ensuremath{^2}(h\ensuremath{-}3s)\ensuremath{-}a\ensuremath{^2}(h\ensuremath{-}3s)\ensuremath{^2}/h=3a\ensuremath{^2}s(h\ensuremath{-}3s)/h. Rearrangement gives 9as\ensuremath{\leq}(2a+1)h. If h>6s, that inequality is immediate from(2a+1)h>6s(2a+1)>9as. This proves the code bound in all cases. The separation of two codewords also gives h\ensuremath{\geq}4s, because one C\_v has at least s members while its weight is at most h\ensuremath{-}3s.

The degree bound gives n\ensuremath{\geq}3s+1\ensuremath{\geq}7; hence a\ensuremath{\geq}ceil((n\ensuremath{-}1)/3)\ensuremath{\geq}2, so two codewords do exist. Combining h\ensuremath{\geq}4s with n\ensuremath{-}1\ensuremath{\leq}3a and n=a+h yields 2n\ensuremath{\geq}12s\ensuremath{-}1 and therefore n\ensuremath{\geq}6s. Now divide the code bound by a>0. The function(2a+1)(n\ensuremath{-}a)/a=2n\ensuremath{-}2a+n/a\ensuremath{-}1 is decreasing in positive a. Substituting the smaller integer b=floor((n+1)/3) gives 9bs\ensuremath{\leq}(2b+1)(n\ensuremath{-}b). Substituting the real lower value(n\ensuremath{-}1)/3 instead gives 27s(n\ensuremath{-}1)\ensuremath{\leq}(2n+1)\ensuremath{^2}.

Finally n\ensuremath{\geq}6s\ensuremath{\geq}12. If27s\ensuremath{\geq}4n+9, the quadratic inequality would force(4n+9)(n\ensuremath{-}1)\ensuremath{\leq}(2n+1)\ensuremath{^2}, or n\ensuremath{\leq}10, a contradiction. Therefore 27s\ensuremath{\leq}4n+8. Replacing s by r+1 yields 27r+19\ensuremath{\leq}4n, as claimed. This derives all three displayed inequalities from the graph hypotheses, without an assumed coding model.

\subsection{What the finite improvement accomplishes}
The located prior explicit bound is n\ensuremath{\geq}5r+6. After integer rounding, the new linear condition ties at r=1 and improves the finite lower requirement for larger r. However, Skottova and Steiner also prove an eventual restriction r=O(n/(log n)\textasciicircum{}c), which is stronger asymptotically. The new theorem does not supersede that result or construct any graph meeting its obstruction.

In particular it does not solve existence for r\ensuremath{\geq}2, cover k>4, or address infinite graphs. Known r=1 constructions are prior evidence, not a new contribution of this packet.

\subsection{Formal evidence and paper form}
The recorded independent isolated Lean 4.34.1 replay compiled the three proof modules and a separate target audit from unchanged frozen sources, with an initially empty own-output directory. Third-party caches were reused and all final endpoints report standard axioms. Fresh replay outcomes belong in the common register.

A short extremal-graph paper should state the three inequalities, explain the criticality-to-code separation, and give small-r comparisons. The precise theorem is more useful than an inflated conjecture-solution headline. Credit the source definitions, prior critical-edge-set paper and classical finite coding inequalities.

{\small This short manuscript records the construction and selected endpoints. The companion Original results anthology contains the extended ordinary proof exposition and its explicitly stated premises; the preserved Lean source remains the complete formal proof record.\par}

\subsection{Verification and reproducibility}
Fresh execution: sakana-erdos944-order-bound: PASS; 3/3 custom modules compiled successfully; receipt Verification/sakana-erdos944-order-bound-fresh-build.json. Compiler pins, source hashes and endpoint axiom outputs are in Verification. Historical receipts and finite checks do not replace fresh replay.

\begin{samepage}
\subsection{Erdos944Bounds.erdos944\_order\_linear}
\begin{quote}\small\ttfamily theorem erdos944\_order\_linear \{r : \ensuremath{\mathbb{N}}\} (hr : 1 \ensuremath{\leq} r)\newline     (hG : G.IsCritical 4)\newline     (hE : \ensuremath{\forall} edges : Set (Sym2 V), G.IsCriticalEdges edges \ensuremath{\to} r < edges.ncard) :\newline     27*r + 19 \ensuremath{\leq} 4 * Fintype.card V\end{quote}
{\small Source: proofs/erdos944-order-bound/OpHack/Erdos944Bounds.lean:42.\par}

{\small Source SHA-256: 8f9e544d0e9f465f8d92d0c0b878ab2bd9b542c0c969eaa4e2d7c908acec27f4\par}

\end{samepage}
\begin{samepage}
\subsection{Erdos944Bounds.erdos944\_order\_quadratic}
\begin{quote}\small\ttfamily theorem erdos944\_order\_quadratic \{r : \ensuremath{\mathbb{N}}\} (hr : 1 \ensuremath{\leq} r)\newline     (hG : G.IsCritical 4)\newline     (hE : \ensuremath{\forall} edges : Set (Sym2 V), G.IsCriticalEdges edges \ensuremath{\to} r < edges.ncard) :\newline     27 * (r+1) * (Fintype.card V - 1) \ensuremath{\leq} (2 * Fintype.card V + 1)\textasciicircum{}2\end{quote}
{\small Source: proofs/erdos944-order-bound/OpHack/Erdos944Bounds.lean:28.\par}

{\small Source SHA-256: 8f9e544d0e9f465f8d92d0c0b878ab2bd9b542c0c969eaa4e2d7c908acec27f4\par}

{\small\textbf{Sources and attribution:} \href{https://www.erdosproblems.com/944}{1. www.erdosproblems.com / 944}; \href{https://arxiv.org/html/2508.08703}{2. arXiv source: 2508.08703}; \href{https://github.com/google-deepmind/formal-conjectures/blob/main/FormalConjectures/ErdosProblems/944.lean}{3. google-deepmind/formal-conjectures / 944.lean}; \href{https://github.com/KitaKen1/erdos-944-dirac-k4-lean}{4. KitaKen1/erdos-944-dirac-k4-lean}; \href{https://arxiv.org/abs/2606.18462}{5. arXiv source: 2606.18462}; \href{https://github.com/alejandrozu/ulam-opdp-difficulty-atlas}{6. alejandrozu/ulam-opdp-difficulty-atlas}; \href{https://github.com/alejandrozu/openmath-2026-judging/blob/d0ed487f487fa7ec814ff705ab55249983fe8707/OpenMath-Judging/review/entries/E02-erdos944-order-bound.txt}{7. alejandrozu/openmath-2026-judging / E02-erdos944-order-bound.txt}; \href{https://github.com/alejandrozu/openmath-2026-judging}{Central source and adjudication archive}.\par}

\end{samepage}
\end{document}
